% Part: history
% Chapter: biographies
% Section: alonzo-church

\documentclass[../../../include/open-logic-section]{subfiles}

\begin{document}

\olfileid{his}{bio}{chu}

\olsection{ألونزو تشرش}

\olphoto{church-alonzo}{Alonzo Church}

وُلد ألونزو تشرش في واشنطن العاصمة في 14 يونيو 1903. وفي طفولته المبكرة
أفقدته حادثة بمسدس هوائي البصر في إحدى عينيه. أنهى المدرسة الإعدادية في
كونيتيكت عام 1920، وبدأ دراسته الجامعية في برينستون في العام نفسه. وأتم
دراسته للدكتوراه عام 1927. وبعد عامين قضاهما في الخارج، عاد تشرش إلى
برينستون. عُرف تشرش بأدبه الشديد ودقته البالغة؛ فكانت كتابته على السبورة
غايةً في الإتقان، وكان يصون الأوراق المهمة بتغطيتها بعناية بإسمنت «دوكو»
(وهو غراء شفاف). وخارج اهتماماته الأكاديمية، كان يستمتع بقراءة مجلات
الخيال العلمي، ولم يكن يتردد في مراسلة المحررين إذا عثر على معلومة غير
دقيقة في ما نُشر.

كانت إنجازات تشرش الأكاديمية عظيمة. فقد طوّر مع طالبيه ستيفن كليني وباركلي
روسر نظريةً في القابلية الفعالة للحساب، هي حساب لامبدا، على نحو مستقل عن
تطوير آلان تورنغ لآلة تورنغ. والتعريفان للقابلية للحساب متكافئان، ومنهما
نشأت ما تُعرف اليوم بـ\emph{أطروحة تشرش--تورنغ}: تكون دالة على الأعداد
الطبيعية قابلةً للحساب الفعال إذا وفقط إذا كانت قابلة للحساب بآلة تورنغ
(أو بحساب لامبدا). وبرهن أيضاً ما يُعرف الآن بـ\emph{مبرهنة تشرش}: مسألة
القرار المتعلقة بصحة صيغ الرتبة الأولى غير قابلة للحل.

واصل تشرش عمله حتى سن متقدمة. وفي عام 1967 غادر برينستون إلى جامعة
كاليفورنيا في لوس أنجلوس، وظل أستاذاً فيها حتى تقاعده عام 1990. وتوفي في
1 أغسطس 1995 عن عمر 92 عاماً.

\begin{reading}
لسيرة موجزة لتشرش، انظر \citet{EndertonND}. أما كتاباته الأصلية في حساب
لامبدا ومسألة القرار (أطروحة تشرش) فهي \citet{Church1936,Church1936a}.
ويسجل \citet{Aspray1984} مقابلةً مع تشرش عن مجتمع الرياضيات في برينستون
في ثلاثينيات القرن العشرين. وكتب تشرش سلسلةً من مراجعات الكتب في
\emph{مجلة المنطق الرمزي} من عام 1936 حتى عام 1979. وهي محفوظة كلها في
موقع جون مكفارلين \citep{MacFarlane2015}.
\end{reading}

\end{document}
